\documentclass[11pt,a4paper]{article}
\usepackage[utf8]{inputenc}
\usepackage[french]{babel}
\usepackage[T1]{fontenc}
\usepackage{geometry}
\usepackage{array}
\usepackage{tabularx}
\usepackage{xcolor}
\usepackage{titlesec}

\geometry{margin=2cm}

\titleformat{\section}{\large\bfseries\color{blue!70!black}}{}{0em}{}[\titlerule]
\titleformat{\subsection}{\normalsize\bfseries}{}{0em}{}

\begin{document}

\begin{center}
    \LARGE \textbf{Administration Réseaux : Guide Complet Linux} \\
    \large Synthèse de cours — Maël \& Collègue — BTS SIO
\end{center}

\vspace{0.5cm}

\section{I. Structure Fondamentale du Shell}
Une commande Bash est interprétée par le système selon une syntaxe précise :
\begin{itemize}
    \item \textbf{Le Programme :} L'exécutable appelé (ex: \texttt{ls}, \texttt{ps}, \texttt{cat}).
    \item \textbf{L'Option :} Modifie le comportement (ex: \texttt{-l}). Les options peuvent être combinées ou interchangées.
    \item \textbf{L'Argument :} La cible (ex: un chemin ou une variable d'environnement comme \texttt{\$HOME}).
\end{itemize}

\section{II. Gestion des Permissions (chmod)}
Les droits d'accès sont répartis entre le propriétaire (\textbf{u}), le groupe (\textbf{g}) et les autres (\textbf{o}).

\subsection{1. Valeurs Octales et Calcul}
Chaque permission correspond à une valeur numérique cumulative :
\begin{itemize}
    \item \textbf{4 (Read - r) :} Autorise la lecture du contenu.
    \item \textbf{2 (Write - w) :} Autorise la modification ou suppression.
    \item \textbf{1 (Execute - x) :} Autorise l'exécution ou l'accès à un dossier.
\end{itemize}

\subsection{2. Sécurité Critique}
Il est impératif de ne jamais accorder simultanément les droits \textbf{w} (2) et \textbf{x} (1) sur un même fichier. Cette configuration permet d'injecter du code malveillant et de l'exécuter immédiatement.

\section{III. Répertoire Exhaustif des Commandes}

\begin{tabularx}{\textwidth}{|l|l|X|}
\hline
\textbf{Commande} & \textbf{Options Clés} & \textbf{Description et Usage} \\
\hline
\texttt{ls} & \texttt{-lah} & Liste le contenu. \texttt{-l} (format long), \texttt{-a} (tous/cachés), \texttt{-h} (taille lisible). \\
\hline
\texttt{cd} & & \textit{Change Directory}. Ramène à la racine utilisateur (\texttt{\$HOME}) par défaut. \\
\hline
\texttt{pwd} & & Affiche le chemin absolu du répertoire de travail actuel. \\
\hline
\texttt{mkdir} & \texttt{-p} & Crée un dossier. \texttt{-p} crée aussi les dossiers parents si nécessaire. \\
\hline
\texttt{touch} & & Crée un fichier vide ou met à jour sa date de modification. \\
\hline
\texttt{cat} & & Affiche le contenu d'un fichier ou permet de concaténer. \\
\hline
\texttt{nano} & & Éditeur de texte simple pour modifier des fichiers en terminal. \\
\hline
\texttt{rm} & \texttt{-rf} & Supprime. \texttt{-r} (récursif), \texttt{-f} (force sans confirmation). \textbf{Attention : } \texttt{rm -rf *}. \\
\hline
\texttt{cp} & \texttt{-R} & Copie des fichiers ou répertoires (mode récursif requis pour les dossiers). \\
\hline
\texttt{mv} & & Déplace ou renomme un fichier/répertoire. \\
\hline
\texttt{ps} & \texttt{-lah} & Visualise l'état des processus en cours d'exécution. \\
\hline
\texttt{clear} & & Nettoie l'écran du terminal et vide le buffer d'affichage. \\
\hline
\end{tabularx}

\section{IV. Administration Système et Réseau}
\begin{itemize}
    \item \textbf{\texttt{sudo}} : \textit{Substitute User Do}. Exécute une commande avec les privilèges root.
    \item \textbf{\texttt{su}} : \textit{Substitute User}. Connexion temporaire (ex: \texttt{su root} ou \texttt{su <user>}).
    \item \textbf{\texttt{scp}} : Copie sécurisée de fichiers entre machines via le protocole SSH.
    \item \textbf{Gestion des comptes :}
    \begin{itemize}
        \item \texttt{useradd} / \texttt{adduser} : Création d'un nouvel utilisateur.
        \item \texttt{usermod} : Modification d'un compte (ex: changement de groupe).
        \item \texttt{groupadd} / \texttt{addgroup} : Création d'un groupe propriétaire.
        \item \texttt{passwd} : Modification du mot de passe (essentiel pour la sécurité).
    \end{itemize}
\end{itemize}

\end{document}
